% why/sections/alignment_is_not_a_curve.tex

\section{Alignment}

\subsection{Motivation}

\begin{remark}
In classical geometry, a railway track is often represented as a curve
\(\gamma(s)\) in space.
\end{remark}

\begin{remark}
This representation is insufficient for engineering purposes, as it
does not capture the intrinsic quantities governing construction,
dynamics, and realization.
\end{remark}

\begin{remark}
In particular, curvature and cant are not derived properties, but
primary design variables.
\end{remark}

\subsection{Definition}

\begin{definition}[Alignment]
An \defterm{alignment} is a structured object defined over arc-length
\(s \in [0,L]\), consisting of the following components:

\begin{itemize}
\item a curvature function
\[
\kappa : [0,L] \rightarrow \mathbb{R},
\]

\item a cant function
\[
\phi : [0,L] \rightarrow \mathbb{R},
\]

\item an initial pose
\[
(\gamma_0, \theta_0),
\]
\end{itemize}

together with the induced state variables
\[
\theta(s) = \theta_0 + \int_0^s \kappa(\sigma)\,\dd\sigma,
\]
\[
\gamma(s) = \gamma_0 + \int_0^s (\cos\theta(\sigma), \sin\theta(\sigma))\,\dd\sigma.
\]
\end{definition}

\subsection{Interpretation}

\begin{remark}
An alignment is not defined by its position \(\gamma(s)\), but by its
curvature \(\kappa(s)\).
\end{remark}

\begin{remark}
The curve \(\gamma(s)\) is a derived quantity obtained by integration.
\end{remark}

\begin{remark}
Thus, curvature is the primary design variable, while position is a
consequence.
\end{remark}

\subsection{State Structure}

\begin{definition}[Alignment state]
The \defterm{state} of an alignment at position \(s\) is given by
\[
(\gamma(s), \theta(s), \kappa(s), \phi(s)).
\]
\end{definition}

\begin{remark}
This state couples geometry, orientation, and dynamics in a unified
representation.
\end{remark}

\subsection{Structural Nature}

\begin{proposition}
An alignment is a low-dimensional structured object, even though its
realization may involve high-dimensional data.
\end{proposition}

\begin{remark}
This structure enables:
\begin{itemize}
\item efficient storage,
\item robust optimization,
\item meaningful interpretation.
\end{itemize}
\end{remark}

\subsection{Relation to Realization}

\begin{remark}
The alignment represents the intended geometry.
\end{remark}

\begin{remark}
The physically constructed track is given by the realization operator:
\[
\gamma_{\mathrm{real}} = \mathcal{R}(\kappa, \phi).
\]
\end{remark}

\begin{remark}
Thus, the alignment exists in design space, while the track exists in
physical space.
\end{remark}

\subsection{Relation to Data}

\begin{remark}
Measured data typically provide sampled positions or polylines.
\end{remark}

\begin{remark}
An alignment must be reconstructed from such data via inverse
problems.
\end{remark}

\subsection{Key Insight}

\begin{proposition}
An alignment is not a curve, but a curvature-driven generative model of
a curve.
\end{proposition}

\subsection{Summary}

\begin{summary}
The alignment is the central object of railway geometry.

It is defined by curvature and cant, from which all geometric and
dynamic quantities are derived.

This perspective shifts the focus from describing curves to defining
generative structures, forming the foundation of
alignment-based information modelling.
\end{summary}
